\section*{Chapter \ref{ch:dv:jumps-and-levy-models} --- Jumps and Lévy Models}

\begin{solution}{exo:dv:jumps-and-levy-models:1}
The number of jumps is Poisson with mean $\lambda T$: $1-e^{-3/52}=5.6\%$ in a week, $1-e^{-3}=95.0\%$ in a year.
\end{solution}

\begin{solution}{exo:dv:jumps-and-levy-models:2}
At $u=-i$, $iu=1$ and $u^2=-1$, so the exponent is
$T\bigl(-\frac12\sigma^2(1-1)-\lambda\bar k+\lambda(e^{\mu_J+\frac12\delta^2}-1)\bigr)=0$. $\phi(-i)=\E[e^{x_T}]=\E[S_T]/F_T$:
the model's forward equals the market's, the martingale condition that the drift correction $-\lambda\bar k$
enforces.
\end{solution}

\begin{solution}{exo:dv:jumps-and-levy-models:3}
Skewness scales as $T^{-1/2}$: $-0.4\sqrt{52}=-2.88$ at one week and $-0.4/\sqrt4=-0.20$ at four years.
\end{solution}

\begin{solution}{exo:dv:jumps-and-levy-models:4}
Hold $\Delta$ shares against the short straddle; the P\&L in each state is $-(V_i-V_0)+\Delta(S_i-S_0)$. Setting
the two equal (a riskless position across the event) gives $\Delta^*=(V_u-V_d)/(S_u-S_d)$. Here
$(2.21+1.48)/(50\times(0.221+0.148))=3.69/18.45=0.20$.
\end{solution}

\begin{solution}{exo:dv:jumps-and-levy-models:5}
The model delta is the slope of the straddle's value for a small move of the spot before the event, where the
straddle is nearly symmetric: $-0.02$. The event is not a small move; the straddle gains on both outcomes, but
more on the up move, so the right hedge is the chord slope $+0.20$. A short position of 0.02 shares adds to the
exposure on the up move instead of reducing it.
\end{solution}

\begin{solution}{exo:dv:jumps-and-levy-models:6}
In a Lévy model the skewness decays like $T^{-1/2}$ and the smile flattens quickly, so a fit at one month leaves
the one-year skew near zero ($-0.07$ against $-0.25$). The market's long skew comes from the persistence and
correlation of volatility, which Bates adds through Heston's variance factor; its jumps keep the short end steep.
\end{solution}

\begin{solution}{exo:dv:jumps-and-levy-models:7}
The standard deviation falls from 0.97 to 0.55 premiums, and the one-in-a-hundred loss from 4.1 to 2.1: smaller,
more frequent jumps are closer to a diffusion. (The total variance also falls, to 13.5\% a year, and the premium
to 1.55.)
\end{solution}

\begin{solution}{exo:dv:jumps-and-levy-models:8}
Most of the 160\% is one jump on Thursday, not a diffusion. Across the jump the Black delta leaves a P\&L standard
deviation of about 11\% of the premium even with known outcome sizes, and a hedge adjusted to the jump leaves 16\%
once the sizes are uncertain. The risk is the size and direction of the move, which only options (or not being
short the event) can hedge.
\end{solution}

\begin{solution}{pb:dv:jumps-and-levy-models:1}
\textbf{1.} 159.5\% at the money, 79.3\% at 38 and 116.4\% at 62.
\textbf{2.} The density is bimodal: two humps around the two outcomes, with little mass exactly at the money. Options
near the money are expensive relative to a lognormal of the same wings, which is a frown.
\textbf{3.} $1.595^2/52-0.35^2/52=0.0465$; a standard deviation of 21.6\%.
\textbf{4.} $21.6\%\times\sqrt{2/\pi}=17.2\%$. The straddle pays the absolute move, so its price measures the expected
absolute move whatever the distribution; the standard deviation is read through a normal assumption that a binary
move violates.
\textbf{5.} The distribution is not lognormal: one volatility matches one strike's price, not the shape.
\textbf{6.} $\sigma=40.6\%$, $p=0.399$, $a=0.199$, with an error of 0.07 volatility point.
\textbf{7.} $+22.1\%$ and $-14.7\%$.
\textbf{8.} The uncertainty of the outcome sizes widens each hump; the two-point model can only widen them through
the diffusion.
\textbf{9.} About 40\%: the market's price of success. A trader compares it with his own estimate (from trial data,
precedents) and trades the difference with digitals or call spreads.
\textbf{10.} $(1-0.4\times1.2214)/0.6-1=-14.8\%$.
\textbf{11.} 8.86; 1.81.
\textbf{12.} $+0.09$ and $-0.02$; 1.01 and 2.01.
\textbf{13.} $0.20$; zero.
\textbf{14.} $0.20$; 1.45, 16.5\% of the premium.
\textbf{15.} Options on the event: buying back part of the straddle, or strangles and butterflies that carry the
exposure to the outcome sizes.
\textbf{16.} Delta is the response to small moves; the ratio is the chord across the two outcomes. They coincide only
for a linear payoff.
\textbf{17.} Delta-hedge the diffusion normally until the event; just before it, move to the chord ratio for the
two outcomes; after it, back to the delta.
\textbf{18.} At least the residual of the best hedge, about 1.45 per straddle (one standard deviation), more for a
confidence level; scaled by the position and aggregated across events.
\textbf{19.} A success probability of 0.40 with moves of $+22.1\%$ and $-14.7\%$; the Black delta leaves a P\&L
standard deviation of 1.01 (11.4\% of the premium) across the event, the binomial ratio 0.20 removes it with known
sizes and leaves 1.45 (16.5\%) with uncertain sizes.
\textbf{20.} The distribution of the move, and with it the hedge ratio across the jump, which no single volatility
contains.
\end{solution}

\begin{solution}{iq:dv:jumps-and-levy-models:1}
In a diffusion a fixed out-of-the-money option needs a move of many standard deviations as $T\to0$, so its price
vanishes exponentially fast; with jumps one jump suffices and the price is proportional to $T$. Market prices of
short-dated out-of-the-money options decline roughly linearly, and short-dated smiles are steep.
\iqlookfor{the speeds of convergence, and the Carr--Wu idea.}
\end{solution}

\begin{solution}{iq:dv:jumps-and-levy-models:2}
$dS/S=-\lambda\bar k\,dt+\sigma\,dW+(e^J-1)\,dN$, $J\sim\mathcal N(\mu_J,\delta^2)$, $N$ Poisson($\lambda$). Given $n$
jumps the log-return is normal, so $C=\sum_n e^{-\lambda T}(\lambda T)^n/n!\;C_{\mathrm{BS}}(F_n,K,T,\sigma_n)$ with
$F_n=Fe^{n(\mu_J+\delta^2/2)-\lambda\bar kT}$ and $\sigma_n^2=\sigma^2+n\delta^2/T$.
\iqlookfor{the compensator, and the conditioning argument.}
\end{solution}

\begin{solution}{iq:dv:jumps-and-levy-models:3}
No: across a jump the hedged position earns $V(Se^J)-V(S)-\Delta S(e^J-1)$, which no single $\Delta$ removes for
all $J$. With two known outcomes a chord ratio does; otherwise hedge the residual convexity with options (they have
the same jump exposure), size positions to the residual, and hold a reserve.
\iqlookfor{the convexity over the jump, and options as the hedge.}
\end{solution}

\begin{solution}{iq:dv:jumps-and-levy-models:4}
Cumulants are linear in $T$, so skewness $\propto T^{-1/2}$ and excess kurtosis $\propto T^{-1}$: steep smiles at
short expiries that flatten quickly. Index smiles flatten more slowly, which calls for stochastic volatility as well.
\iqlookfor{the scaling, and its consequence.}
\end{solution}

\begin{solution}{iq:dv:jumps-and-levy-models:5}
The variance gamma characteristic function decays only like $|u|^{-2T/\nu}$, very slowly at one week, while Heston's
decays exponentially; a grid truncated at a moderate $u$ loses a large tail. Extend the grid (or use a damped or
tail-corrected scheme), and test against the conditional closed form at the shortest expiry.
\iqlookfor{the decay rate of the characteristic function.}
\end{solution}

\begin{solution}{iq:dv:jumps-and-levy-models:6}
Separate the diffusion from the event: price with a diffusion at the ordinary volatility plus a jump whose law is read
from the smile (an \omterm{def:dv:parametrising-the-surface:event}{event variance}, or a two-point fit if the smile shows a frown). Hedge the diffusion with delta
before the event and the jump with the chord ratio or with options; after the event, implied volatility collapses to
the diffusion. What can go wrong: the size or direction of the move, a move larger than priced, liquidity at the open,
and a crush of the other expiries' volatility.
\iqlookfor{diffusion plus event, the hedge across the jump, the volatility crush.}
\end{solution}
